\chapterimage{figs/chapterhead.png}
\chapter{Parser, expressões e linguagem interna}
\label{chap:parser_expressoes_linguagem_interna}

\section{Introdução}
\label{sec:cap9_introducao}

Ao longo da primeira metade deste livro, a linguagem do \textit{GenESyS} apareceu para o usuário principalmente por meio da aba \textit{Simulang}. Naquela etapa, o interesse central estava em perceber que o modelo gráfico também possui uma representação textual. Para o desenvolvedor, porém, essa constatação é apenas o ponto de partida. O problema real agora é outro: como o sistema interpreta expressões, como a linguagem se liga ao estado do modelo e como novas capacidades sintáticas e semânticas podem ser introduzidas no simulador.

Este capítulo trata exatamente desse subsistema. Seu foco não está apenas na noção genérica de ``parser'', mas no modo concreto como o \textit{GenESyS} organiza a avaliação de expressões, o tratamento de símbolos do modelo, a integração com funções estocásticas, a gramática Bison/Flex e os pontos de extensão que permitem ao sistema crescer.

A importância desse tema é grande. Em um simulador como o \textit{GenESyS}, a linguagem não é ornamento. Muitas decisões do modelo passam por expressões: parâmetros, fórmulas, referências a atributos, consultas a estatísticas, controles de simulação e funções probabilísticas. O parser, portanto, não é apenas uma conveniência textual; ele é parte da semântica operacional do simulador.

\section{O papel do parser no GenESyS}
\label{sec:cap9_papel_parser_genesys}

O parser do \textit{GenESyS} deve ser entendido como a infraestrutura que torna possível interpretar expressões e funções em contexto de modelo. Isso significa que ele não opera apenas sobre texto abstrato, mas sobre texto que precisa ser relacionado a elementos concretos do simulador:
\begin{itemize}
    \item atributos da entidade corrente;
    \item controles e respostas de simulação;
    \item contadores e coletores estatísticos;
    \item variáveis, fórmulas, filas, recursos e outros objetos do modelo;
    \item funções matemáticas e probabilísticas;
    \item informações sobre o estado corrente da simulação.
\end{itemize}

Em outras palavras, o parser funciona como ponte entre linguagem e estado interno do sistema. Essa é a razão pela qual ele é tão importante na arquitetura do GenESyS.

\subsection{Parser como subsistema semântico}
\label{subsec:cap9_parser_como_subsistema_semantico}

É útil evitar uma leitura limitada do parser como simples analisador sintático. No \textit{GenESyS}, ele faz mais do que reconhecer uma expressão bem formada. Ele também precisa:
\begin{itemize}
    \item resolver nomes de símbolos;
    \item identificar a natureza desses símbolos;
    \item obter valores associados ao estado atual do modelo;
    \item integrar funções estocásticas apoiadas em um amostrador;
    \item propagar mensagens de erro úteis.
\end{itemize}

Isso significa que a leitura correta do parser no GenESyS é semântica e não apenas sintática.

\section{A interface \texttt{Parser\_if}}
\label{sec:cap9_interface_parser_if}

A interface \texttt{Parser\_if} oferece uma síntese muito clara do papel do parser no projeto. Ela define um contrato mínimo para avaliação de expressões e funções estocásticas, deixando explícitas algumas responsabilidades centrais:
\begin{itemize}
    \item avaliar expressões e retornar um valor numérico;
    \item oferecer uma variante que reporte sucesso e mensagem de erro sem lançar exceção;
    \item disponibilizar a última mensagem de erro;
    \item associar o parser a uma instância de \texttt{Sampler\_if};
    \item expor o driver interno do parser para cenários avançados.
\end{itemize}

Essa interface é importante porque mostra que o parser é visto, arquiteturalmente, como serviço do simulador e não como detalhe acidental de implementação.

\subsection{Avaliação de expressões}
\label{subsec:cap9_avaliacao_expressoes}

O ponto mais evidente da interface é o método de avaliação de expressões. A existência de duas formas principais de \texttt{parse} é particularmente reveladora:
\begin{itemize}
    \item uma forma que pode lançar exceções;
    \item uma forma que evita exceções e informa explicitamente sucesso e mensagem de erro.
\end{itemize}

Essa decisão melhora a integração do parser com diferentes camadas do sistema. Em alguns contextos, o desenvolvedor pode preferir tratamento por exceção. Em outros, especialmente em validação e interface de usuário, é mais conveniente receber o estado de sucesso de forma explícita.

\subsection{Integração com \texttt{Sampler\_if}}
\label{subsec:cap9_integracao_sampler_if}

A interface também deixa explícito que o parser não é apenas analisador de expressões determinísticas. Ele integra funções estocásticas apoiadas em \texttt{Sampler\_if}. Isso é importante porque reforça a natureza específica da linguagem do GenESyS: ela não é apenas uma mini-linguagem aritmética, mas uma linguagem voltada à simulação.

Em termos práticos, isso significa que funções probabilísticas da linguagem dependem da presença de um amostrador associado ao parser. Com isso, a semântica das expressões incorpora naturalmente mecanismos de geração aleatória.

\section{A classe \texttt{Model} como usuária do parser}
\label{sec:cap9_model_como_usuario_parser}

O papel do parser dentro do kernel fica especialmente claro na classe \texttt{Model}. Ela expõe métodos como \texttt{parseExpression} e \texttt{checkExpression}, o que mostra que o modelo usa o parser diretamente para:
\begin{itemize}
    \item avaliar expressões definidas no contexto do modelo;
    \item validar expressões inseridas pelo usuário;
    \item verificar referências a \textit{data definitions} no texto das expressões.
\end{itemize}

Esse detalhe é arquiteturalmente muito importante. Ele significa que o parser não está isolado em uma camada externa; ele participa da vida do modelo.

\subsection{Por que \texttt{Model} invoca o parser}
\label{subsec:cap9_por_que_model_invoca_parser}

A presença desses métodos em \texttt{Model} revela um princípio importante: expressões fazem parte da semântica do modelo, e não apenas da interface. Um componente ou dado do modelo que armazena uma expressão precisa poder validá-la e, em muitos casos, avaliá-la dentro do contexto corrente da simulação.

Essa escolha arquitetural é muito boa porque mantém a linguagem conectada ao objeto ao qual ela realmente pertence: o modelo.

\subsection{Expressões e verificação de consistência}
\label{subsec:cap9_expressoes_verificacao_consistencia}

Também é importante notar que a classe \texttt{Model} oferece \texttt{checkExpression}. Isso mostra que o parser participa não apenas da execução, mas também da fase de verificação do modelo. Em outras palavras, a validade de uma expressão é tratada como parte da integridade do modelo, e não como problema a ser descoberto apenas durante uma execução já iniciada.

Essa decisão melhora muito a robustez do sistema e ajuda a explicar por que a verificação do modelo é uma etapa tão importante no uso e no desenvolvimento do GenESyS.

\section{A gramática Bison/Flex do GenESyS}
\label{sec:cap9_gramatica_bison_flex}

A implementação concreta do parser em \texttt{source/parser/parserBisonFlex} mostra de forma bastante clara a maturidade desse subsistema. O arquivo \texttt{bisonparser.yy} define uma gramática C++ baseada em \texttt{lalr1.cc}, com classe de parser própria, tokens tipados, localizações, código de suporte e integração explícita com o \texttt{genesyspp\_driver}. Isso mostra que o parser foi pensado como componente sério da infraestrutura do simulador, e não como rotina textual improvisada.

A gramática cobre, entre outros elementos:
\begin{itemize}
    \item números;
    \item operadores relacionais e lógicos;
    \item funções trigonométricas e matemáticas;
    \item funções probabilísticas;
    \item funções associadas ao kernel;
    \item símbolos do modelo;
    \item extensões introduzidas por \textit{plugins}.
\end{itemize}

Essa abrangência já basta para mostrar que o parser é um dos subsistemas de maior densidade semântica do projeto.

\subsection{Uma linguagem realmente integrada ao modelo}
\label{subsec:cap9_linguagem_integrada_modelo}

A leitura da gramática mostra algo particularmente relevante: muitos tokens e produções não se referem apenas a construções sintáticas abstratas, mas a objetos reais do modelo. Há regras para atributos, respostas de simulação, controles de simulação e elementos como variáveis, fórmulas, filas e recursos. Também aparecem funções associadas a estatísticas, contadores e estado da simulação.

Isso reforça um ponto essencial:
\begin{quote}
    A linguagem do GenESyS não é externa ao simulador; ela é uma projeção textual da semântica do modelo e da simulação.
\end{quote}

\subsection{Funções probabilísticas e amostragem}
\label{subsec:cap9_funcoes_probabilisticas_amostragem}

A gramática deixa explícita a presença de funções probabilísticas como exponencial, normal, uniforme, Weibull, lognormal, gama, Erlang, triangular e beta, entre outras. Em todas elas, a semântica operacional depende do \texttt{Sampler\_if} associado ao parser.

Isso é muito importante para o desenvolvedor, porque mostra que:
\begin{itemize}
    \item a linguagem do simulador incorpora diretamente amostragem estocástica;
    \item a implementação do parser precisa manter contrato consistente com o subsistema estatístico;
    \item o tratamento de erros nessas funções não é apenas sintático, mas também operacional.
\end{itemize}

\section{Símbolos do modelo e dependência do estado corrente}
\label{sec:cap9_simbolos_modelo_estado_corrente}

Uma das características mais interessantes do parser do GenESyS é sua dependência explícita do estado corrente da simulação. A gramática mostra, por exemplo, que o acesso a atributos depende da entidade corrente associada ao evento atual. Da mesma forma, funções como \texttt{TNOW}, \texttt{TFIN}, número de repetições e identificadores da entidade corrente dependem do estado mantido em \texttt{ModelSimulation} e no evento atual.

Essa característica é arquiteturalmente forte porque revela que a interpretação de expressões no GenESyS é contextual. A mesma expressão pode ter significado diferente de acordo com:
\begin{itemize}
    \item o modelo carregado;
    \item a entidade corrente;
    \item o evento em processamento;
    \item o instante simulado;
    \item os dados do modelo disponíveis naquele contexto.
\end{itemize}

\subsection{Parser como consumidor do estado do kernel}
\label{subsec:cap9_parser_consumidor_estado_kernel}

Isso significa que o parser consome diretamente informações do kernel. Ele consulta:
\begin{itemize}
    \item o modelo;
    \item a simulação;
    \item o evento corrente;
    \item o \textit{data manager};
    \item coletores e contadores;
    \item controles e respostas.
\end{itemize}

Essa dependência não é defeito. Pelo contrário, é parte da natureza do problema. Um parser para linguagem de simulação precisa conhecer o universo de símbolos e estados do simulador.

\section{Pontos de extensão por \textit{plugins}}
\label{sec:cap9_pontos_extensao_plugins}

A gramática do \textit{GenESyS} revela de forma muito clara a existência de regiões pensadas para extensão. Os blocos marcados por comentários estruturados para inclusões, tokens, tipos e produções de \textit{plugins} mostram que o sistema foi desenhado para crescer.

Esse aspecto é particularmente importante para o desenvolvedor. Ele mostra que a ampliação da linguagem não depende necessariamente de reescrita completa da gramática base. Há pontos arquiteturalmente previstos para que novos elementos sejam introduzidos.

\subsection{Extensão sintática e semântica}
\label{subsec:cap9_extensao_sintatica_semantica}

Quando um \textit{plugin} amplia a linguagem, ele não introduz apenas novos tokens. Em geral, ele introduz também:
\begin{itemize}
    \item novos símbolos;
    \item novas funções;
    \item novas regras de interpretação;
    \item novas relações entre texto e dados do modelo.
\end{itemize}

Isso significa que a extensão do parser é simultaneamente sintática e semântica. É precisamente por isso que a infraestrutura de \texttt{ParserChangesInformation} e \texttt{ParserManager} existe no projeto.

\section{A classe \texttt{ParserManager}}
\label{sec:cap9_classe_parsermanager}

A classe \texttt{ParserManager} torna explícita a ambição arquitetural do GenESyS de sustentar evolução da linguagem. Sua interface prevê:
\begin{itemize}
    \item geração de um novo parser a partir de mudanças;
    \item conexão de um novo parser ao sistema.
\end{itemize}

Embora o próprio cabeçalho indique que partes importantes desse fluxo ainda estejam em desenvolvimento, a existência dessa classe já é arquiteturalmente relevante. Ela mostra que o GenESyS não trata a gramática como artefato inteiramente estático. Há, no desenho do sistema, espaço para gestão explícita da evolução do parser.

\subsection{O que isso significa para o desenvolvedor}
\label{subsec:cap9_o_que_significa_para_desenvolvedor}

Para o desenvolvedor, isso significa que trabalhar com linguagem no GenESyS não se limita a modificar um arquivo Bison ou Flex. O projeto prevê, em nível arquitetural, um subsistema responsável por:
\begin{itemize}
    \item descrever mudanças;
    \item gerar artefatos de parser;
    \item conectar novos parsers ao simulador.
\end{itemize}

Mesmo que essa infraestrutura ainda esteja amadurecendo, ela já orienta a forma correta de pensar a evolução da linguagem.

\section{Erros, mensagens e robustez}
\label{sec:cap9_erros_mensagens_robustez}

Outro aspecto importante da implementação do parser é o tratamento de erros. A gramática mostra preocupação explícita com mensagens úteis, inclusive no caso de tokens ilegais, literais não encontrados e falhas em funções probabilísticas. Isso é importante porque a linguagem de simulação, quando mal interpretada ou mal diagnosticada, se torna rapidamente fonte de frustração para usuário e desenvolvedor.

A robustez do parser depende, portanto, de três fatores:
\begin{itemize}
    \item reconhecimento sintático correto;
    \item integração semântica correta com o modelo;
    \item produção de mensagens de erro suficientemente informativas.
\end{itemize}

Essa tríade é especialmente importante em um sistema em que expressões podem aparecer em muitos pontos do modelo.

\section{Como estudar e modificar o parser de forma produtiva}
\label{sec:cap9_como_estudar_modificar_parser}

A forma mais produtiva de estudar o parser do GenESyS é em camadas.

Primeiro, compreender a interface \texttt{Parser\_if} e a forma como \texttt{Model} invoca o parser. Depois, estudar a gramática Bison/Flex e a maneira como ela acessa símbolos do modelo. Em seguida, identificar os pontos de extensão introduzidos por \textit{plugins}. Só então faz sentido propor alterações mais significativas, especialmente quando envolvem novas funções, novos símbolos ou mudanças na gramática.

Também é importante lembrar que alterações no parser não são apenas alterações textuais. Elas podem afetar:
\begin{itemize}
    \item verificação de modelos;
    \item semântica de componentes e \textit{data definitions};
    \item integração com amostragem;
    \item mensagens de erro;
    \item compatibilidade com \textit{plugins}.
\end{itemize}

\begin{table}[htbp]
    \centering
    \caption{Camadas de leitura do parser do GenESyS.}
    \label{tab:cap9_camadas_leitura_parser}
    \small
    \begin{tabular}{|p{3.9cm}|p{7.3cm}|}
        \hline
        \textbf{Camada} & \textbf{Pergunta principal} \\
        \hline
        Interface & Que serviço o parser oferece ao restante do sistema? \\
        \hline
        Integração com o modelo & Como \texttt{Model} e o estado da simulação utilizam esse serviço? \\
        \hline
        Gramática concreta & Como a linguagem é reconhecida e avaliada? \\
        \hline
        Extensibilidade & Como novos símbolos e funções entram no sistema? \\
        \hline
        Robustez & Como erros são detectados e comunicados? \\
        \hline
    \end{tabular}
\end{table}

\section{Considerações Finais}
\label{sec:cap9_consideracoes_finais}

Este capítulo mostrou que o parser do \textit{GenESyS} é muito mais do que uma camada auxiliar de texto. Ele é parte integrante da semântica do simulador. A interface \texttt{Parser\_if}, a integração direta com \texttt{Model}, a gramática Bison/Flex, o uso de símbolos do modelo e da simulação, a dependência de um amostrador e os pontos explícitos de extensão por \textit{plugins} revelam um subsistema arquiteturalmente denso e estrategicamente importante para o projeto.

Com isso, o leitor já dispõe dos elementos principais para compreender como o \textit{GenESyS} liga linguagem, modelo e execução. O passo seguinte é ampliar novamente o foco, saindo do parser e voltando ao ecossistema mais amplo do projeto: aplicações, ferramentas, exemplos em C++, testes e estratégias de evolução do software. É exatamente isso que será tratado no próximo capítulo, que fecha esta versão do manual do desenvolvedor.

Teste seu entendimento desse conteúdo respondendo as seguintes perguntas:

\begin{enumerate}
    \item Por que o parser do \textit{GenESyS} deve ser entendido como subsistema semântico e não apenas sintático?

    \item Qual é a importância de a classe \texttt{Model} expor métodos como \texttt{parseExpression} e \texttt{checkExpression}?

    \item Como a gramática do \textit{GenESyS} evidencia dependência do estado corrente do modelo e da simulação?

    \item Por que os pontos de extensão por \textit{plugins} são essenciais para a evolução da linguagem do simulador?
\end{enumerate}